\chapter{Método e falsificação}
\label{chap:metodo}

Benchmarking de sistemas é sensível a fatores aparentemente inócuos
\parencite{mytkowicz2009wrongdata}; por isso, protocolo, ordem, tolerâncias e
critérios de rejeição devem ser congelados antes dos resultados. A análise
estatística deve respeitar pareamento e custo de repetição
\parencite{kalibera2013rigorous}.

\section{Decisões críticas}

\begin{table}[htbp]
\caption{Superfícies mínimas de autoridade causal.}
\label{tab:decisoes}
\centering
\begin{tabularx}{\textwidth}{@{}rXll@{}}
\toprule
\# & Decisão & Consumo & Teste causal \\
\midrule
1 & Identidade e integridade do dataset & antes do load & troca de hash \\
2 & Bundle de inicialização & antes dos parâmetros & seed adulterada \\
3 & Precisão numérica & antes da primitive & mudança de dtype \\
4 & Função de perda & antes do forward & loss alternativa \\
5 & Regra de atualização & antes do step & learning rate \\
6 & Estado do otimizador & em cada transição & estado reaproveitado \\
7 & Working set & antes do acesso & subconjunto distinto \\
8 & Parada & antes de continuar & critério distinto \\
9 & Equivalência e aceitação & antes da publicação & tolerância distinta \\
\bottomrule
\end{tabularx}
\end{table}

\section{Exemplo de política}

\begin{lstlisting}[language=TLang,caption={Fragmento ilustrativo; substituir por
fonte T real e hashada no manuscrito final.},label={lst:policy}]
policy acceptance(receipt, evidence) {
  require evidence.dataset_hash == receipt.dataset_hash;
  require evidence.optimizer_steps == 400;
  if evidence.max_parameter_difference <= 1e-9 {
    accept("operational-equivalence");
  } else {
    reject("semantic-divergence");
  }
}
\end{lstlisting}
